\chapter{Практика}
\label{ch:chap2}

Выполним реализацию ФВЧ с параметрами:

\begin{enumerate}
    \item Подавление в полосе задерживания А = 53db
    \item Частота пропускания $f_{pass}$ = 2000 Hz
    \item Частота непропускания $f_{stop}$ = 500 Hz
    \item Переходная полоса $\triangle f = f_{pass} - f_{stop}$
    \item Частота пропускания $f_c = \frac{f_{pass} + f_{stop}}{2}$ 
    \item Частота дискретизации $f_s$ = 30000Гц
\end{enumerate}

\section{ИХ ФВЧ}

\begin{figure}[H]
    \centering
    \includegraphics[width=1.0\textwidth]{IR.png}
    \caption{ИХ ФВЧ}
\end{figure}

Здесь $\Omega_c$ - нормированная (цифровая частота), которая равна $\frac{2\pi f}{f_s}$. Измеряется в $\frac{rad}{sample}$. \\

Изобразим ИХ, предварительно рассчитав кол-во необходимых отсчетов ИХ по эмпирической формуле
$N = 2M+1$, где $N = \frac{3.3}{\triangle f}$, где $\triangle f = \frac{f_{pass} - f_{stop}}{f_s}$. Отсюда можем 
найти M как N//2. Также изобразим сдвинутую ИХ (строим от 0 до 2M, а не от -M до M).

\begin{figure}[H]
    \centering
    \includegraphics[width=1.0\textwidth]{IR_py.png}
    \caption{Визуализация ИХ ФВЧ}
\end{figure}

\section{Частотная характеристика ФВЧ}
Выполним ППФ над ИХ и получим частотную характеристику ФВЧ. Дополнительно изобразим частотную характеристику в логарифмическом 
масштабе (db):

\begin{figure}[H]
    \centering
    \includegraphics[width=1.0\textwidth]{db.png}
    \caption{Визуализация частотной характеристики}
\end{figure}

В нашем случае $f_c = 1750Hz$, а на графике при -3db получаем значение $\approx 1800Hz$, что соответствует $f_c$.

\begin{figure}[H]
    \centering
    \includegraphics[width=1.0\textwidth]{lin.png}
    \caption{Визуализация частотной характеристики}
\end{figure}

\section{Тестовый сигнал}
Сгенерируем сигнал, состоящий из двух компонент: $f_1 = 4000Hz$, $f_2 = 1200Hz$. Изобразим его спектр

\begin{figure}[H]
    \centering
    \includegraphics[width=1.0\textwidth]{signal.png}
    \caption{Тестовый сигнал в частотной и временной области}
\end{figure}

\section{Фильтрация тестового сигнала}
С помощью ФВЧ отфильтруем сигнал и изобразим спектр полученного сигнала. В идеале спектральная линия на 1200Hz должна исчезнуть.

\begin{figure}[H]
    \centering
    \includegraphics[width=1.0\textwidth]{spec.png}
    \caption{Спектр отфильтрованного сигнала}
\end{figure}

Линия не исчезла до конца, небольшое кол-во энергии осталось на частоте 1200Hz. 

\section{Весовая функция}
Для решения проблемы выше используются весовые функции или окна (Ханна, Хемминга т.п). Это окно перемножается с ИХ фильтра,
чтобы сгладить боковые лепестки. Изобразим окно Хемминга, ИХ и результат их перемножения:

\begin{figure}[H]
    \centering
    \includegraphics[width=1.0\textwidth]{win.png}
    \caption{Окно Хемминга и ИХ ФВЧ}
\end{figure}

Можно заметить, что "хвосты" ИХ сгладились. \\

Теперь пропустим сигнал через новую ИХ:

\begin{figure}[H]
    \centering
    \includegraphics[width=1.0\textwidth]{spec2.png}
    \caption{Сигнал на выходе фильтра}
\end{figure}

Можем заметить, что проблема решена - спектр сигнала стал "чистым".


\endinput 